\scanpage{166}

AK

\uline{Pseudomorphismes}

$X$ un préschéma. ouvert schématiquement dense $U$ : si $X$ est le plus \deleted{petit} ss-préschéma fermé de $X$ dominant $U$ i.e. si $\overline U=X$ [lorsque $U\longrightarrow X$ est quasi-cpt cela implique que pour tout ouvert $V$, $U\cap V$ est schématiquement dense dans $X$ -- aussi cela vaut-il mieux comme définition : toute \deleted{fonction}\inserted{section de $\underline O_X$} sur un ouvert $V$ de $X$ nulle sur $U\cap V$ est nulle]

% slanted
\margin{\unsure{condition équiv. à ci-dessous}}

\deleted{\unsure{L'intéresse|L'intérêt}} Caractère local de la notion [Si $X=\bigcup X_i$, les $X_i$ ouverts, alors $U$, schématiquement dense ds $X\Longleftrightarrow$ les $U\cap X_i$ sch. dense dans les $X_i$]

$U$ schémati. dense dans $X\Rightarrow U$ schém. dense dans \deleted{des} $V$ \unsure{q.c.q.|q.p.|q.c.} $U\subset V\subset X$

Si $U,U'$ sont schém. denses, leur $U\cap U'$ l'est

[Soit $X$ et $Y$ sur $S$, $\deleted{X}Y$ séparé sur $S$, $\deleted{X}\unsure{\Rightarrow|\longrightarrow}U$ schém. dense dans $X$, si $f,g:X\longrightarrow Y$ sont des $S$-morphismes tels que $f\mid U=g\mid U$

\deleted{\unsure{Quasi}}\inserted{Pseudo}-morphismes : \unsure{classes|classe} d'équivalence d'\deleted{applications}morphismes définis sur des ouverts schématiquement denses, comme $S$-morphismes. Pour un quasi-$S$-morphisme dans un $S$-\deleted{pré}\uline{schéma}, il y a un \inserted{\uline{(plus grand)}} domaine de définition.

\[
\deleted{\text{Le domaine }D(f\mid V)={}}\quad
\mathrm{Dom}(f\mid V)=\mathrm{Dom}(f)\cap V
\]

% archive sheet number: 165
